\clearpage
\phantomsection
\addcontentsline{toc}{section}{Resumen}

\noindent{\bfseries\large Resumen} \\[0.5cm]

El presente proyecto de innovación y mejora tuvo como objetivo principal desarrollar e implementar un sistema web con flujos automatizados 
para la gestión de nóminas y procesos de recursos humanos orientado a los clientes empresariales de J\&P~Periféricos~S.A.C. 
En el diagnóstico de la situación inicial, se evidenció que el procesamiento mensual de planillas en las organizaciones atendidas dependía de hojas de cálculo dispersas, 
consolidación manual de asistencias y validaciones repetitivas, lo que ocasionaba demoras prolongadas en los cierres de pago, errores en el cálculo de horas extras, 
bonificaciones y retenciones de ley (ONP, AFP y quinta categoría), así como riesgos de incumplimiento normativo y filtración de información financiera sensible.

Para resolver esta problemática, el proyecto se ejecutó bajo el marco de trabajo ágil Scrum, 
estructurando el desarrollo en iteraciones incrementales alineadas con los requerimientos funcionales y no funcionales del negocio. 
A nivel tecnológico, se diseñó una arquitectura monolítica modular desacoplada, empleando TypeScript, NestJS y Fastify en el backend, 
React en el frontend y PostgreSQL como motor de base de datos relacional normalizado. Asimismo, se integró la plataforma de orquestación \textbf{n8n} 
para ejecutar flujos de trabajo asíncronos encargados de generar automáticamente los archivos planos y estructuras regulatorias exigidas 
por SUNAT PLAME, AFPnet y las entidades bancarias (telecrédito), además de emitir y distribuir por correo electrónico las boletas de pago 
en formato PDF protegidas mediante cifrado individual con el número de documento de identidad de cada trabajador. En materia de seguridad de la información, 
el sistema implementó control de acceso basado en roles (RBAC), cifrado simétrico \textbf{AES-256-GCM} 
para atributos financieros críticos y un módulo de auditoría inmutable para garantizar la trazabilidad integral de las operaciones.

En cuanto a su organización,
 el documento se estructura en ocho capítulos: 
 el Capítulo~1 describe las generalidades de J\&P~Periféricos~S.A.C. y el levantamiento inicial mediante entrevistas; 
 el Capítulo~2 plantea la problemática, los objetivos, los antecedentes, la justificación y el marco teórico-normativo; 
 el Capítulo~3 analiza la situación actual (\textit{AS-IS}) mediante diagramas de procesos, Ishikawa y priorización de Pareto; 
 el Capítulo~4 detalla el plan de acción bajo Scrum, los requerimientos, la arquitectura, los procesos mejorados (\textit{TO-BE}) y el cronograma; 
 los Capítulos~5 y 6 desarrollan la estructura de costos y la evaluación técnico-económica; 
 finalmente, los Capítulos~7 y 8 presentan los resultados, la discusión, las conclusiones y las recomendaciones para la puesta en producción y evolución futura del producto.